% ===================================================================
% CHAPTER 1: INTRODUCTION
% ===================================================================
\chapter{Introduction}
\label{ch:intro}

\section{Adaptive Navigation in Vector Field Environments}
\label{sec:ch1:motivation}
% Ported from: thesis.tex Ch.1 Motivation (lines 332-346); extended with
% cwaight_systems_paper.tex Sec. I (convergence zones, leaks, eddies) and
% Draft_11.tex Sec. I (coherent structures, transport barriers).

Many environments of scientific and practical interest are organized by an
underlying vector field: an ocean current, an atmospheric flow, a diffusing
plume, a magnetic field. Locating and characterizing the special points and
structures of such a field, its sinks, sources, vortices, saddles, and the
separatrices that divide one flow regime from another, is often more useful
than simply following the field's local direction. A search-and-rescue team
drifting with an ocean current benefits from knowing where a nearby eddy
will trap debris; an environmental monitoring team benefits from knowing
where two water masses meet. In many of these settings no map of the field
exists in advance, and no single sensor can see the whole field at once. A
team of mobile robots, each sampling the field locally and sharing
measurements with its neighbors, can estimate field structure that no
individual robot could determine alone, and can then use that estimate to
steer.

The oceanic and atmospheric environments that motivate this work are
governed by a spatially varying velocity or force, and their transport and
accumulation behavior is organized by the critical points of the field, the
coordinates where the flow vanishes \cite{1,2}. A surface
convergence zone acts as a sink, gathering buoyant debris and drifting
bodies; this makes such zones high probability search regions after a loss
at sea \cite{3,4}. Containing a ruptured wellhead or a harbor
discharge requires reaching the source against its own outward flow
\cite{5,6}. A mesoscale ocean eddy serves as a center, trapping a
body of water, along with the plankton bloom inside it, and carrying both
across a basin for weeks \cite{7,8,9}. Monitoring the bloom
over its lifetime therefore requires a sensor to orbit the drifting core at
a fixed standoff; a ground-fixed station would not stay with the eddy.

Beyond critical points, vector fields have curves that are useful for
navigation. Dynamic ocean flows contain naturally occurring boundaries known
as coherent structures, the strain-dominated curves across which water
masses separate and along which floating material collects
\cite{10,11}. Mapping these curves lets a team deploy resources to a
targeted search area instead of sampling a much larger basin to find where
the structure lies \cite{11,12,13,14}.

The gap this dissertation addresses is locating and characterizing these
features, points and curves alike, from measurements a robot team takes of
the physical field itself, not from an artificial field constructed to
guide the team, and not from a heading alone. A heading toward, or along, a
feature is not the same as knowing where the feature is.

\section{Literature Review}
\label{sec:ch1:literature}

\subsection{Adaptive Navigation in Scalar Fields}
\label{sec:ch1:lit-scalar}
% Ported from: Draft_11.tex Sec. I (lines 94-106), extended with
% cwaight_systems_paper.tex Sec. I (line 51) for the strategy list and
% field-deployment citations.

Adaptive navigation (AN) in scalar fields is well developed. Each robot
contributes one scalar measurement to a policy that reaches and follows a
feature, and formations are sized so that differential measurements across
the cluster recover the gradient and curvature that policy consumes
\cite{15,16}. \"{O}gren, Fiorelli, and Leonard established cooperative
gradient climbing with a formation of mobile sensors \cite{16}.
Bri\~{n}\'{o}n-Arranz, Renzaglia, and Schenato extended it to the Hessian,
five robots on a ring plus one at the center recovering both in closed form
\cite{15}. Kitts, McDonald, and Neumann developed primitives for extrema
finding, contour following, ridge/trench following, and saddle point
station keeping \cite{17}, and verified ridge, trench, saddle, and front
tracking experimentally on an indoor testbed \cite{18,19}.

Additional strategies extend this line of work to gradient-based cluster
space navigation \cite{20}, extended information consensus filters
\cite{21,22}, constrained nonsmooth distributed optimization
\cite{23}, and bio-inspired search \cite{24}, and several such
systems have been field-deployed on ground, surface, and aerial vehicles for
environmental monitoring, pollution tracking, and search and rescue
\cite{20,14,12,25}.

\subsection{Artificial Vector Fields for Guidance}
\label{sec:ch1:lit-artificial}
% Ported from: cwaight_systems_paper.tex Sec. I (line 53), expanded with
% individual citations, including Liddy et al. from thesis.tex, naming the
% examples the source sentence only cited as a group.

Much of the broader vector field literature constructs artificial fields
for path-following and obstacle avoidance, a separate problem from sensing
a real flow, and field campaigns that track features carried by an actual
ocean current show the difference \cite{26}. Representative examples
include vector field path following for miniature air vehicles
\cite{27}, singularity-free guiding vector fields for robot navigation
\cite{28}, obstacle avoidance using complex vector fields
\cite{29}, motion planning and collision avoidance using navigation
vector fields \cite{30}, and adaptive vector field guidance without a
priori knowledge of course dynamics and wind \cite{31}. None of these
fields are sensed; they are specified in advance, and the robot's task is
to follow one, a different problem from the one this dissertation poses.

\subsection{Robotic Sensing of Physical Flows}
\label{sec:ch1:lit-sensing}
% Ported from: Draft_11.tex Sec. I (lines 116-124), extended with
% cwaight_systems_paper.tex Sec. I (lines 53, 57) for the two-camps framing
% and the Knizhnik citation.

Work that senses the physical field directly falls into two camps. One
camp presumes prior knowledge of the environment, a map, a model, or an
assumed functional form, unavailable in the unmapped settings this
dissertation targets \cite{32,33}. The other operates online using
only local measurements. Michini et al. bracket a presumed stable or
unstable manifold with three robots using only local velocity measurements
\cite{34}, an approach inspired by the PIM triple procedure, extended
to $N$ robots \cite{35} and validated experimentally in gyre-like flows
\cite{36}. Knizhnik et al. address a related setting, flow-based
control of marine robots in gyre-like environments \cite{37}. In the
Michini line of work the robots begin on the manifold and are advected by
the flow, and the manifold's identity is supplied in advance; the saddle
point itself is treated by neither the controller nor its analysis, and
\cite{34} reports veering away from the structure near a saddle in
flow-tank experiments. Kularatne and Hsieh confirm the boundary by
computing the local finite-time Lyapunov exponent field on board, at the
cost of a nonzero integration horizon and a commitment to attracting
structures \cite{1}. In every case the robotic formation pre-straddles
the structure and does not acquire a structure it does not already occupy.

\section{Problem Statement}
\label{sec:ch1:problem}
% Ported from: thesis.tex Ch.1 Problem Statement (lines 365-388).

Existing approaches to multirobot navigation in scalar and vector fields
generally fall into two categories. The first requires prior knowledge of
the field, a map, a model, or an assumed functional form, against which a
robot team's local measurements are compared or fit. The second uses only
local information but provides directional guidance alone: climb the
gradient, follow the current, avoid the strain. Neither category directly
answers the question this dissertation poses: given only local, real-time
measurements and no prior model, can a robot team estimate the location and
type of a field's critical points, and, more generally, the field's local
first- and second-order structure, well enough to navigate with respect to
it, not just along it.

This dissertation answers that question constructively. It shows that
three robots are both necessary and sufficient to recover a
two-dimensional vector field's local linearization (its Jacobian) from
simultaneous local measurements alone, and that six robots are both
necessary and (almost always) sufficient to recover the same field's local
quadratic structure (its Jacobian, Okubo-Weiss parameter, and Hessian). In
both cases the estimate is exact for the class of fields the robots'
formation size implies (affine for three robots, quadratic for six),
assembled without any assumption about the field's global shape, and used
directly inside a navigation control law whose stability is established
analytically and validated in simulation and, where hardware permits, on a
physical testbed.

\section{Thesis Statement: Fit Order Sets the Reachable Feature Class}
\label{sec:ch1:thesis}
% New text. States the fit-order ladder, cites cwaight_systems_paper.tex
% Sec. VI-B (vector-sum/vector-to-scalar drift) and Draft_11.tex Sec. I/VII
% (linear-fit-at-one-point limitation, cubic future work) as the motivating
% failures each rung fixes.

This dissertation is organized around one idea: the order of the local
polynomial a robot team fits to a sensed vector field sets which class of
field feature that team can locate. A zeroth-order fit returns a single
vector, the field's local mean. A first-order (affine) fit returns the
field's local Jacobian. A second-order (quadratic) fit returns the field's
local Hessian in addition to the Jacobian. A third-order (cubic) fit would
return the field's local third derivatives. Each order needs more robots,
sampled at one instant, to determine without assuming information the
field has not supplied, and each order unlocks a feature class the order
below cannot represent. Table~\ref{tab:ch1:fit-order-ladder} states the
ladder used throughout this dissertation.

\begin{table}[h]
\centering
\caption{The fit-order ladder used throughout this dissertation.}
\label{tab:ch1:fit-order-ladder}
\small
\begin{tabular}{@{}p{1.3in}cp{1.1in}p{1.7in}c@{}}
\toprule
Fit order & Robots per component & Degeneracy & Features reachable & Chapter \\
\midrule
Zeroth (mean vector, magnitude) & 1+ & none & heading only; drifts on orbits & \ref{ch:zeroth} \\
First (affine, Jacobian) & 3 & collinear robots & critical points: location, type, orbit & \ref{ch:first_order} \\
Second (quadratic) & 6 & robots on a common conic & separatrices, Okubo-Weiss boundaries, saddles as minima & \ref{ch:second_order} \\
Third (cubic, future) & 10 & robots on a common cubic & true Hessian of $D$, transverse curvature & \ref{ch:limitations} \\
\bottomrule
\end{tabular}
\end{table}

Chapter~\ref{ch:zeroth} shows why the ladder starts here. A zeroth-order
primitive steers on the mean of the perceived flow direction across the
formation (the vector-sum primitive) or on the flow's magnitude alone (the
vector-to-scalar primitive \cite{38}, here modified to move tangentially
along the direction of steepest descent). Both supply a heading, not a
location. Orbiting a vortex under either primitive drifts outward in
spiral fashion, because neither has a mechanism to detect or correct
radial error with respect to a specific point, and the vector-sum drift has
been confirmed in hardware \cite{39}. On a saddle the failure is worse:
the field vectors do not all point toward the critical point, so the
vector-sum controller cannot produce an orbital trajectory at all.

A first-order fit removes this failure. Recovering the field's Jacobian at
a single point gives both the location and the type of an isolated critical
point, and supports a control law that converges onto it or holds a
commanded orbit around it (Chapter~\ref{ch:first_order}). But a linear fit
is a property of one point \cite{40}. It cannot by itself locate an
extended curve, such as the separatrix or Okubo-Weiss boundary that
organizes transport in a region no single critical point describes, because
a one-point fit has no gradient of its own to climb. A second-order fit
gives the velocity gradient as a field over the formation rather than a
value at a point, so a quantity built from it, the Jacobian determinant or
the smaller rate-of-strain eigenvalue, comes with its own gradient in
closed form and can be climbed to and held on a curve
(Chapter~\ref{ch:second_order}).

The next rung is not built in this dissertation. Chapter~\ref{ch:limitations}
describes a cubic fit from ten robots as future work: it would recover the
field's third derivatives directly, giving the true Hessian of the
determinant $D$ and the transverse curvature of a rate-of-strain trench,
quantities the second-order fit can only approximate.

\section{Contributions}
\label{sec:ch1:contributions}
% Ported from: thesis.tex Ch.1 Contributions (lines 389-424), rewritten in
% fit-order terms. The old list's second item, a compositional grammar for
% general cluster-space kinematics, is not listed here; it is out of this
% dissertation's outline, which derives only the three-robot and six-robot
% kinematics used by the first- and second-order chapters (Appendices
% \ref{app:kin3} and \ref{app:kin6}).

The main contributions of this dissertation are:
\begin{enumerate}
  \item A hardware and simulation toolkit for multirobot vector-field
  navigation: HSV-encoded printed fields and neural-network sensor
  calibration and recalibration (R\textsuperscript{2}~=~0.96 direction,
  0.91 magnitude) \cite{39}, robot dynamics identification, and a
  matching simulator (Chapter~\ref{ch:tools}). The Decabot platform itself
  is prior work of other members of the Robotic Systems Laboratory; the
  contribution here is the sensing, calibration, and simulation layer built
  on it.
  \item Zeroth-order baseline primitives, vector-sum and vector-to-scalar
  direction-only navigation from prior work \cite{38}, implemented and
  evaluated on the hardware testbed of Chapter~\ref{ch:tools}
  \cite{39}, establishing what direction
  alone can and cannot certify (Chapter~\ref{ch:zeroth}).
  \item A distributed, model-free method for three robots to estimate a
  two-dimensional vector field's local Jacobian from simultaneous local
  measurements alone, with no prior field knowledge, proven minimal at
  exactly three robots, enabling closed-form critical point localization,
  eigenvalue-based type classification, and attraction and orbital control
  laws; validated in simulation (eight fields, 1000 Monte Carlo trials each,
  100\% convergence) and on hardware (157 convergence trials plus 12
  orbital trials, 100\% success, sub-centimeter bias)
  (Chapter~\ref{ch:first_order}).
  \item A distributed, model-free method for six robots (pentagon plus
  center) to estimate a two-dimensional flow's local Jacobian, Okubo-Weiss
  parameter, and Hessian from simultaneous local measurements, proven
  minimal at exactly six robots and almost-always sufficient, ruling out
  the degenerate hexagon and conic trap, enabling two modified-Newton
  controllers, the $D$ tracker and the $s_1$ tracker, that track the
  Okubo-Weiss boundary and the separatrix trench respectively, with a
  proven stability and traversal guarantee and validation on both a
  canonical double-gyre model and real ocean high-frequency-radar current
  data (Chapter~\ref{ch:second_order}).
  \item Findings that hold across fit orders: how estimation noise
  amplifies with fit order, how formation size trades truncation error
  against noise sensitivity, and how objectivity under a rotating observer
  depends on which surrogate field a controller climbs
  (Chapter~\ref{ch:crosscutting}).
\end{enumerate}

\section{Publications}
\label{sec:ch1:publications}
% Ported from: thesis.tex Ch.1 Publications (lines 426-453), restructured
% to the four publications behind this dissertation's fit-order chapters.
% Titles taken verbatim from each source's own title. Statuses left as
% [TBD] where no source states one; per CLAUDE.md, this dissertation's own
% submission is intentionally held until the critical-points paper
% publishes, and that is not a blocker to flag.

Portions of this dissertation have appeared, or are in preparation to
appear, in the following publications.

\begin{table}[h]
\centering
\caption{Mapping of dissertation chapters to publications.}
\label{tab:ch1:publications}
\small
\begin{tabular}{@{}cp{1.9in}p{1.0in}c@{}}
\toprule
Chapter & Working title & Venue & Status \\
\midrule
\ref{ch:tools}, \ref{ch:zeroth} & A Functional Indoor Testbed for Multirobot Adaptive
Navigation in Vector Field Environments \cite{39} & ASME IDETC/CIE
2025, Anaheim (DETC2025-167604) & Published \\
\ref{ch:first_order} & Adaptive Navigation of Multirobot Systems to
Critical Points in 2D Vector Fields in Simulation and Experiment & IEEE
Systems Journal & \textbf{[TBD]} \\
\ref{ch:second_order} & Second-Order Cooperative Field Estimation for
Multirobot Tracking of Coherent Flow Structures & IEEE Systems Journal &
In preparation \\
\bottomrule
\end{tabular}
\end{table}
%% IDETC 2025 testbed paper: published (author confirmed 2026-09-29; final
%% PDF in Paper_Writing/Reference Papers/[33]).
%% The 2024 IDETC submission (DETC2024-142746) was rejected; its early
%% ideas are not part of this dissertation and it is not cited.
%% TODO(status): confirm the critical-points paper's submission status.

\section{Dissertation Organization}
\label{sec:ch1:organization}
% Ported from: thesis.tex Ch.1 Organization (lines 454-474), rewritten for
% the ten-chapter, six-appendix fit-order structure.

Chapter~\ref{ch:tools} describes the verification tools this dissertation
depends on: the Decabot hardware platform (prior work of the Robotic
Systems Laboratory), the HSV-encoded printed vector fields and
neural-network sensor calibration developed for it, which converts a
camera reading to a velocity estimate \cite{39}, robot dynamics
identification, and the matching simulation environment used throughout
the rest of this dissertation.

Chapter~\ref{ch:preliminaries} sets out the mathematics common to every fit
order: the planar vector field, the velocity gradient tensor, critical
point classification by eigenvalue structure, the rotation-strain
decomposition and the Okubo-Weiss partition, coherent structures,
objectivity under a change of observer frame, and the canonical benchmark
fields used for validation.

Chapter~\ref{ch:estimation} presents the cooperative estimation framework
shared by every fit order above zeroth: the multilayer control
architecture, cluster space control for the three-robot and six-robot
formations, local polynomial field estimation and its minimality,
estimator sensitivity to noise and truncation, and the fit-order ladder of
Section~\ref{sec:ch1:thesis} restated with this framework's machinery.

Chapter~\ref{ch:zeroth} presents the zeroth-order baseline: direction-only
navigation and the vector-sum and vector-to-scalar primitives \cite{38},
hardware trials on
fixed and sinking vortices on the testbed of Chapter~\ref{ch:tools}
\cite{39}, and the argument that direction alone is not location.

Chapter~\ref{ch:first_order} presents the first-order result: three robots
recovering a vector field's local Jacobian, and the resulting critical
point estimation, classification, attraction, and orbital control laws,
validated in simulation and on hardware and compared against the
zeroth-order primitives of Chapter~\ref{ch:zeroth}.

Chapter~\ref{ch:second_order} presents the second-order result: six robots
recovering a flow's local Hessian and Okubo-Weiss structure, the $D$
tracker and $s_1$ tracker built on that estimate, their stability and
traversal properties, and their validation on a double-gyre benchmark and
on Santa Barbara Channel radar current data.

Chapter~\ref{ch:crosscutting} collects findings that appear only once
results from more than one fit order are set side by side: how noise
amplifies with fit order, how formation size trades truncation error
against noise sensitivity, formation conditioning versus field
conditioning, instantaneous estimation on a time-varying field, and
objectivity and sign references across controllers.

Chapter~\ref{ch:limitations} states this dissertation's limitations by fit
order and describes future work, including the cubic, ten-robot extension
that would complete the next rung of the ladder.

Chapter~\ref{ch:conclusion} summarizes the contributions and closes with
final thoughts.

Appendix~\ref{app:fields} collects the canonical field definitions used in
simulation throughout this dissertation. Appendix~\ref{app:kin3} derives
the three-robot cluster's kinematics in full. Appendix~\ref{app:kin6}
derives the six-robot pentagon-plus-center cluster's kinematics in full.
Appendix~\ref{app:frame_equivariance} proves the frame equivariance of the
$s_1$ tracker. Appendix~\ref{app:stability} collects the transverse
stability proofs for the controllers of Chapters~\ref{ch:first_order}
and~\ref{ch:second_order}. Appendix~\ref{app:stats} documents the
statistical methods used to report simulation and hardware results.

\section{Notation}
\label{sec:ch1:notation}
% Ported from: thesis.tex Ch.1 "Notation and Definitions" (lines 476-586),
% chapter cross-references updated to the ten-chapter structure and
% trimmed of notation tied to material not in this dissertation's outline
% (the consensus orientation gauge and the composition-tree spine angle,
% both part of the compositional-grammar work excluded per Section
% \ref{sec:ch1:contributions}). The full symbol reference is the Glossary
% of Terms in the front matter.

Throughout this dissertation, bold lowercase letters denote vectors, e.g.,
$\mathbf{p} \in \mathbb{R}^2$, and bold uppercase letters denote matrices,
e.g., $\mathbf{J} \in \mathbb{R}^{m \times n}$. The studies collected in
this dissertation were originally written as independent papers, and their
notations were never cross-checked against one another. Most symbols
already agree; a small number collide and are resolved here, once, for the
whole document. Chapter-specific notation not covered here is introduced
where it is first used.

\subsection*{The bold-\texorpdfstring{$\mathbf{p}$}{p} versus scalar-\texorpdfstring{$p$}{p} rule}

Bold $\mathbf{p}$ always denotes a two-dimensional position,
$\mathbf{p} = (x,y)$, and is never written unbolded when used as a
position. Scalar $p$ always denotes the side-length formation parameter
inherited from the Kitts-lab cluster-space lineage, introduced in
Chapter~\ref{ch:estimation} and used through Appendices~\ref{app:kin3} and
\ref{app:kin6}: the distance from robot 1 to robot 2 in a side-angle-side
(SAS) triangular formation. There is no unbolded $p$ for position and no
bold $\mathbf{p}$ for side length. This rule is carried unchanged from the
source papers and holds throughout every chapter of this dissertation.

\subsection*{The bold-\texorpdfstring{$\mathbf{q}$}{q} versus scalar-\texorpdfstring{$q$}{q} rule}

Chapter~\ref{ch:estimation} introduces bold $\mathbf{q}$ as the general
cluster-space state vector (pose and shape stacked, dimension $2N$ for $N$
robots), while scalar $q$, following the same inherited SAS convention as
scalar $p$, denotes the triangle's second side length (the 2--3 distance).
These two uses never appear in the same equation and are distinguished
exactly as bold $\mathbf{p}$ and scalar $p$ are: there is no unbolded
$\mathbf{q}$ for the state vector and no bold $q$ for the side length.

\subsection*{The alpha collision}

The symbol $\alpha$ is used for three unrelated quantities across the
source materials. Each is disambiguated with a subscript in this
dissertation and must never be written as bare $\alpha$ after this point.

\begin{itemize}
\item $\alpha_{\text{eig}}$: the real part of a complex eigenvalue pair of
  the field Jacobian $\mathbf{J}$, written $\alpha_{\text{eig}} \pm
  i\omega$. This is a property of the vector field's linearization at a
  critical point and classifies spiral-type critical points
  (Chapter~\ref{ch:first_order}).
\item $\alpha_{\text{mom}}$: the discrete first-order momentum filter
  coefficient for robot dynamics, $\alpha_{\text{mom}} = \exp(-\Delta
  t/\tau)$. At 10~Hz with $\tau \approx 0.28$~s, $\alpha_{\text{mom}} \approx
  0.70$. This is characterized in Chapter~\ref{ch:tools} and appears
  identically in the robot dynamics of Chapters~\ref{ch:first_order} and
  \ref{ch:second_order}.
\item $\phi_1, \phi_3$: the two SAS-triangle vertex angles at robots 1 and
  3, used only in the triangle's law-of-cosines derivation
  (Appendix~\ref{app:kin3}). These were originally written as $\alpha$ and
  $\gamma$ in one source paper's appendix; they are renamed here since they
  are otherwise unused and the original choice collided with both
  $\alpha_{\text{eig}}$ and $\alpha_{\text{mom}}$.
\end{itemize}

\subsection*{The \texorpdfstring{$r$}{r} collision}

Three meanings of $r$ appear across the source materials: the orbital
radial vector/magnitude $\mathbf{r} = \mathbf{p}^* - \mathbf{p}_c$
(Chapter~\ref{ch:first_order}), the SAS triangle's third side length (the
1--3 diagonal), and the polar radius used in several canonical field
definitions. This dissertation keeps $r$ for the orbital radial quantity,
since it appears in a named control law. The triangle's third side is
renamed $s$, and the polar radius in field definitions is renamed $\rho$,
matching the pentagon ring radius already used in
Chapter~\ref{ch:second_order} and Appendix~\ref{app:kin6}, so that
``radius of a circular construction'' reads consistently throughout.

\subsection*{The \texorpdfstring{$\mathbf{J}$}{J} collision}

Bold $\mathbf{J}$ denotes the field Jacobian, $\mathbf{J} =
\left[\begin{smallmatrix} u_x & u_y \\ v_x & v_y \end{smallmatrix}\right]$,
throughout this dissertation: it is defined in
Chapter~\ref{ch:preliminaries}, recovered from measurements in
Chapter~\ref{ch:estimation}, and is the central object of
Chapter~\ref{ch:first_order} (first order) and Chapter~\ref{ch:second_order}
(embedded in the second-order estimator). Every kinematic or cluster-space
Jacobian, mapping shape-parameter velocities to individual robot
velocities, is written $\mathbf{J}_c$ instead, whether it is the $6\times
6$ triangle Jacobian or the $12\times 12$ pentagon Jacobian, both presented
in Chapter~\ref{ch:estimation} and derived in full in
Appendices~\ref{app:kin3} and~\ref{app:kin6}.

\subsection*{Other symbols}

The full symbol reference for field and estimation notation, critical
point classification, control laws, robot dynamics, and formation
parameters is the Glossary of Terms in the front matter, extended with the
Chapter~\ref{ch:estimation} and Chapter~\ref{ch:second_order} additions
listed there. Of particular note: capital $\boldsymbol{\Phi}$ (the
six-robot formation matrix, Chapter~\ref{ch:second_order}) is visually
close to lowercase $\phi_1, \phi_3$ (the triangle vertex angles,
Appendix~\ref{app:kin3}) in some fonts; the two never appear in the same
chapter, but the reader should not assume a relationship between them.
Similarly, $\rho$ appears in two guises across this dissertation, the field
definitions' polar radius (Chapter~\ref{ch:preliminaries},
Appendix~\ref{app:fields}) and the pentagon formation's ring radius
(Chapter~\ref{ch:second_order}, Appendix~\ref{app:kin6}).

As in Chapter~\ref{ch:first_order}, this dissertation writes a critical
point of a flow, including a saddle point of the Okubo-Weiss field in
Chapter~\ref{ch:second_order}, as $\mathbf{p}^*$.
